\section{The Reserve as an Extractable Subsidy}\label{sec:subsidy}

The reserve of \cref{cr:buffer} makes above-floor loans payable. It does not make them worth repaying. For exit routing, a holder whose tokens would redeem for a gross value $v$ per token receives $(1-f_s)v$ by redeeming and $(1-\phi)a$ by borrowing at the ratio $a$ and never repaying. Abandonment pays whenever $(1-\phi)\,a>(1-f_s)\,v$, the condition under which the floor-cap contract invites abandonment, with $\gamma\Pf$ replaced by $a$ \citep{extended}. For a holder who redeems at the floor, $v=\Pf$, and a ratio $a>\Pf$ satisfies the condition unless the fees are at least as large as the lending premium, $(1-f_s)/(1-\phi)\ge a/\Pf$. Under the floor cap, $a\le\gbar\Pf$, and a cap with $\gbar+\fsbar<1$ gives the opposite inequality at every sell fee within $\fsbar$; lending above the floor removes that guarantee.

Borrowers can do better than choose between repaying in full and abandoning. They can repay in part, redeem the collateral that the repayment releases, and abandon the rest. The next result says exactly how much this takes out.

\begin{proposition}[What a committed reserve can lose]\label{cr:extraction}
In the setting of \cref{cr:rights}, with no fees, let the book be safe under its repayment rights, that is, $W\ge W_{\min}$. Let a coalition hold all $L=S$ tokens, locked against the loans of the book, and let the reserve $W$ have been paid in by someone else.
\begin{enumerate}[label=(\alph*)]
\item If the coalition repays down to an attainable residual configuration with total lock $Q$ and total debt $B$, redeems the released tokens and never repays the rest, it receives $B-R(Q)$ more than by repaying every loan and redeeming every token. The supremum of this gain over attainable configurations, which is the maximum over their closure, is $W_{\min}$, the least reserve that \cref{cr:kernel} requires. Under proportional release the supremum is attained.
\item After a path that attains this maximum the issuer holds cash $W-W_{\min}$.
\item Under proportional release, a book that is admitted with total debt $D_{\max}$ of \eqref{cr:dmax} and $W>0$ has $W_{\min}=W$: the coalition can take the whole reserve.
\item For a coalition that buys the $L$ tokens from an empty curve and pays a buy fee $f_b$, an origination fee $\phi$ and a sell fee $f_s$, all proportional and paid out of the mechanism, the profit of the path in (a) is
\[
\Pi=B-R(Q)-f_b\,R(L)-\phi\,D-f_s\,\bigl[R(L)-R(Q)\bigr].
\]
\item Under the curve-valued contract of \cref{cr:nonlinear} every attainable configuration has $B\le R(Q)$, so the gain in (a) is never positive.
\end{enumerate}
\end{proposition}
\begin{proof}
Because the book is safe under its repayment rights, every repayment and redemption below can be executed; without that hypothesis $B-R(Q)$ is still the difference of the two cash flows as written, but the issuer may run out of cash before the redemption is paid. (a) Repaying down to the configuration costs $D-B$ and the released $L-Q$ tokens redeem for $R(L)-R(Q)$; repaying everything costs $D$ and returns $R(L)$. The difference is $B-R(Q)$. By \cref{cr:kernel} the least admissible reserve is the supremum of $\sum_ib_i-R(\sum_iq_i)$ over attainable configurations, equivalently its maximum over their closure, and it is at least zero because $(0,0)$ is attainable for every loan. (b) The proof of \cref{cr:kernel} gives the cash after the path as $W+R(Q)-B$. (c) Because $W>0$, the function $(R(q)+W)/q$ tends to infinity as $q\downarrow0$ and is continuous on $(0,L]$, so the infimum in \eqref{cr:dmax} is attained at some $q^\ast>0$, where $(D_{\max}/L)\,q^\ast=R(q^\ast)+W$. By concavity $H(q^\ast)\ge D_{\max}q^\ast/L$, so $H(q^\ast)-R(q^\ast)\ge W$, and admission requires $W_{\min}\le W$. (d) The purchase costs $(1+f_b)R(L)$, the loans pay out $(1-\phi)D$, the repayment costs $D-B$ and the redemption returns $(1-f_s)[R(L)-R(Q)]$; the sum is $\Pi$. (e) The persistent condition \eqref{cr:nonlinearcover} holds at every attainable configuration, so $B\le\gbar\sum_iR(q_i)\le R(Q)$ by \eqref{cr:superadd}.
\end{proof}

The reserve that makes a proportional book safe is thus exactly the amount that a coalition of its borrowers can extract from it. Take $R(q)=q+q^2/2$, $L=S=1$ and $W=0.08$, which is below the threshold $kL^2/8=0.125$ at which $D_{\max}$ passes the redemption value $R(L)$ of the whole collateral. Then $D_{\max}=1.4<R(1)=1.5$, so repaying in full is still better than abandoning in full. But buying the token for $1.5$, borrowing $1.4$, repaying $0.84$ to release $0.6$ tokens, redeeming them for $1.02$ and leaving $0.4$ tokens locked against a debt of $0.56$ nets $0.08$: the whole reserve. The issuer is left with no cash and a residual loan of $0.56$ against collateral that redeems for $R(0.4)=0.48$: the loan is under-collateralized, in a state that remains safe under its repayment rights. With all three fees at $1\%$ the profit is $0.0408$ (script block N16). The statement is about a strategy that is available, under the stated assumptions on access and on who funded the reserve; it is not a prediction of what borrowers do, and a borrower who paid in the reserve counts that payment as a cost. The borrower's position is the residual claim of the corporate-debt model, the right to recover the collateral by paying the debt or to walk away \citep{merton1974}, and $W_{\min}$ is capital held against the worst admissible scenario, in the spirit of, though not shown to satisfy the axioms of, a coherent risk measure \citep{artzner1999}.

Two consequences follow. At the maximum admitted credit \eqref{cr:dmax}, the committed reserve is exposed to complete loss under the strategy above: it is first-loss capital for that path, though other repayment paths return it. And the amount that must stay committed is the current $W_{\min}$, not $W$: a repayment narrows the set of residual configurations and can lower the requirement, after which the difference may be withdrawn while the loan is open. In the example, repaying $1.12$ leaves $0.2$ tokens against a debt of $0.28$ and a requirement of $0.06$, so $0.02$ of the reserve is free. A path that repays and abandons can instead use up the cash and leave the remaining requirement committed for good. The curve-valued contract of the next section does not have the defect, by part~(e), and under $\gbar+\fsbar<1$ redeeming a lock pays more than abandoning it.

